Los requisitos del SCDEE se elicitaron a partir de tres fuentes principales:
las reuniones periódicas con el tutor, el estudio comparativo de las
plataformas descritas en \cref{SEC:SOTA-PLATFORMS}, y los requisitos
institucionales transversales de seguridad y privacidad. Cada requisito se
formula operativamente en términos de entrada, proceso y salida, lo que
facilita su validación posterior.

Desde la perspectiva del backend, el actor inmediato es el \dfn{frontend},
que materializa toda interacción con los actores humanos finales mediante
peticiones a la \ac{api} \ac{rest}. Los actores humanos del sistema,
ordenados por nivel de privilegio decreciente, son: el
\textbf{superadministrador} del sistema, que gestiona las organizaciones
cliente; el \textbf{gestor de organización} (\texttt{is\_staff=True}), que
administra usuarios, cursos y configuración dentro de una
organización; el \textbf{coordinador de asignatura}, responsable académico
de una asignatura, con permisos plenos sobre sus exámenes; el
\textbf{profesor}, corrector de las instancias que tiene asignadas; y el
\textbf{alumno}, que consulta su instancia y solicita revisión. El sistema
no implementa una jerarquía rígida de roles globales: los permisos
académicos se definen contextualmente por asignatura mediante
\texttt{SubjectMembership}.
